% The plan grows a second step: chapter 8 answered one query with one fetch;
% here a second subgraph joins, and the same shape of client request now
% needs two fetches, the second of which cannot start until the first has
% come back.
%
% Same idiom as figures/tikz/ch01-one-field.tex and ch08-one-fetch.tex: each
% lane is a timeline read left to right, ticks mark stages, and a brace under
% the stages that matter carries a one-line label. Axis, tick, event, span
% and spanlabel styles are copied from those files unchanged.
%
% Lane A (sessions { nodes { title } }, answered inside one subgraph) and
% lane B (the same field plus speaker { name }, the first query this book
% answers across two subgraphs) share the same first four stages: client
% posts to the router, the router plans, the router posts to Sessions, and
% Sessions runs one statement. Those four ticks sit at the same x in both
% lanes for exactly that reason, so the eye reads the two rows as the same
% request up to the point where they part company. From there lane A is
% finished, so its axis simply stops; lane B keeps going for two more ticks,
% because that is the whole argument of the figure, and a lane that stopped
% at the same width as lane A would be hiding it. This is a deliberate
% widening, not a timing claim: SPEC decision 19 puts no milliseconds in
% this book, the horizontal axis is sequence only, and no gap here is sized
% to suggest one stage took longer than another.
%
% The dashed vertical rule is the seam idiom fixed by
% figures/tikz/ch05-the-join-moves.tex, reused unchanged: a dashed vertical
% marks where work crosses from one subgraph's process into another's. It
% appears once in lane B, between the statement that finishes the Sessions
% fetch and the tick where the router sends ids onward as representations,
% because that is the one crossing this request makes. Lane A crosses no
% seam and carries no dashed rule at all, which is half the figure's point:
% one lane has a seam to cross, the other does not.
%
% Two braces sit under lane B, with a visible gap on either side of the
% seam so neither brace is mistaken for spanning it, the same discipline
% ch08-one-fetch.tex uses for its own pair of braces. The shorter brace, under
% the fetch to Sessions, notes that the router rewrote the selection to ask
% for the id it will need rather than the name the client asked for. The
% main brace, under the fetch to Speakers, names the step lane A never
% takes at all: a fetch that depends on the first one returning, and that
% costs one more statement rather than one per session.
%
% No quotation marks anywhere in this file. The book's strict Ascii mode
% (see check-chapter.psd1) makes the prose gate read a literal " as a quote
% character, while \" is LaTeX's diaeresis accent and fails to compile
% inside a node. Every identifier below is written bare, and URL paths are
% broken onto their own line by hand rather than quoted or wrapped.
%
% Bare tikzpicture (house style): the figure environment, caption and label
% live at the call site in the section file.
\begin{tikzpicture}[
  x=1mm, y=1mm,
  every node/.append style={font=\sffamily\scriptsize},
  axis/.style={draw, line width=0.4pt,
    -{Straight Barb[length=1.4mm, width=1.6mm]}},
  tick/.style={draw, line width=0.4pt},
  event/.style={anchor=south, align=center, text width=20mm,
    inner sep=1pt},
  span/.style={draw, line width=0.4pt},
  spanlabel/.style={anchor=north, align=center, inner sep=2pt},
  lane/.style={anchor=west, font=\sffamily\scriptsize\itshape},
  boundary/.style={draw, line width=0.4pt, dashed},
]

% ------------------------------------------------------------- lane A
% At y=20, level with lane B's own label. The three-line event at x=36
% reaches y=13, so the 13 that ch08-one-fetch.tex uses for its lane A
% would sit on top of it here; that lane's tall labels start further
% right and its title has ended by then.
\node[lane] at (0,20) {A. sessions only: one query, one subgraph};

\draw[axis] (0,0) -- (90,0);
\foreach \x in {14,36,58,82}{
  \draw[tick] (\x,-1.5) -- (\x,1.5);
}

\node[event, text width=20mm] at (14,3)
  {client posts to\\localhost:3002\\/graphql};
\node[event, text width=22mm] at (36,3)
  {router plans:\\one fetch,\\subgraph sessions};
\node[event, text width=22mm] at (58,3)
  {router posts to\\localhost:5001\\/graphql};
\node[event, text width=18mm] at (82,3)
  {1 statement\\to SQLite};

% ------------------------------------------------------------- lane B
% Lane A carries no braces, so it needs less room beneath it than lane B
% does. 35mm puts the two lanes the same distance apart as the two in
% ch08-one-fetch.tex measured from axis to axis.
\begin{scope}[shift={(0,-35)}]
  \node[lane] at (0,20)
    {B. speaker name added: the first query split across two subgraphs};

  \draw[axis] (0,0) -- (150,0);
  \foreach \x in {14,36,58,82,112,140}{
    \draw[tick] (\x,-1.5) -- (\x,1.5);
  }

  \draw[boundary] (97,-2) -- (97,8);
  \node[spanlabel, inner sep=1pt] at (97,-2) {seam};

  \node[event, text width=20mm] at (14,3)
    {client posts to\\localhost:3002\\/graphql};
  \node[event, text width=20mm] at (36,3)
    {router plans:\\Sequence,\\two steps};
  \node[event, text width=22mm] at (58,3)
    {router posts to\\localhost:5001\\/graphql:\\id, not name};
  \node[event, text width=24mm] at (82,3)
    {1 statement\\to SQLite,\\Sessions service};
  \node[event, text width=30mm] at (112,3)
    {router sends ids\\as representations\\to localhost:5002\\/graphql};
  \node[event, text width=24mm] at (140,3)
    {1 statement\\to SQLite,\\Speakers service};

  \draw[span] (58,-4) -- (58,-7) -- (82,-7) -- (82,-4);
  \node[spanlabel] at (70,-8)
    {the router rewrote the selection:\\it asks Sessions for the id it\\will
     need, not the name asked for};

  \draw[span] (112,-4) -- (112,-7) -- (140,-7) -- (140,-4);
  \node[spanlabel] at (126,-8)
    {the step lane A never takes:\\a BatchEntity fetch that waits\\on the
     first fetch, and costs one\\more statement, not one per session};
\end{scope}

\end{tikzpicture}
